% part1_captions.tex -- GENERATED by scripts/figs_part1.py (do not edit by hand; rerun the script).
% Caption macros for figs/{thmA_design,base23_lift,projlines}.pdf; every number is asserted by the script.
\newcommand{\thmADesignCaption}{%
Theorem~\ref{thm:A} at its threshold: the net $\mathcal P^{U^*_4}_{1,q}$ of the lattice design $U^*_4$ ($a^*_4=5$) at $q=5a^*_4+4=29$, $k=1$, $N=q^4$; in cell units (side $N^{-1/4}$). (a) The elementary slab of the cells $(n_3,n_4)=(11,20)$, a $15\times15$ window of cells $(n_1,n_2)$, projected to $(x_1,x_2)$; one point per cell. Open circles: the lattice image $n+r/q$, $r=F_0n\bmod q$ ((H1) of Lemma~\ref{lem:Ss}); dots: the points; grey segments: the jitter of the two deeper digits, inside $[0,1/q)^4$. Accent arrows: the projections of $(I+F_0/q)e_1$ and $(I+F_0/q)e_2$; dashed: slips, where $r_1$ or $r_2$ wraps; crosses: lattice sites the net skips there. (b) Zoom on one slip of $r_1$ (dotted box in (a)), with the jitter boxes of side $1/q$: across the slip the step is $1.828$ instead of $|(I+F_0/q)e_1|$. Proved: $\delta N^{1/4}\ge\mathrm{LB}=24/29=0.8275$, the relaxation of Lemma~\ref{lem:Ss} at $k=1$ (here equal to the bound $1-a^*_4/q$ of Theorem~\ref{thm:A}). Measured (computed, exact): $c=0.8287$; the lattice minimum is $\lambda_1(I+F_0/q)=0.8297$.}
\newcommand{\baseTwoThreeCaption}{%
The deep-shift lift (Definition~\ref{def:lift}) in base $b=3$: the digitally shifted (Definition~\ref{def:dshift}) $(0,9,3)$-net of Theorem~\ref{thm:base23}(1) for $m\equiv0\pmod 3$, one lift of the $m=6$ seed. A point has one top digit per coordinate, $6$ block digits (a copy of a seed point shifted below block depth $\lfloor 6/2\rfloor=3$) and two jitter digits; the top digits $(\alpha,\alpha+\beta_1,\alpha+\beta_1+\beta_2)$ select one of $27$ top boxes (side $\frac13$). (a) The slab $z\in[\frac{11}{27},\frac{12}{27})$ projected to $(x,y)$: 729 points, one in each box of side $3^{-3}$; the nine top boxes (thick) carry the same frozen pattern. (b) The top box $(\alpha,\beta_1,\beta_2)=(1,1,2)$ (dashed in (a)): each point (dot) lies in the frozen box of side $3^{-4}$ (shaded; open circle at its corner) of exactly one point of the seed copy (plus signs), i.e. the box fixed by the top digits of that point, which later lifts never change; grey segments: the jitter, i.e.\ the deep block digits and the two jitter digits. (c) The nine digits of the marked point of (b). Computed, exact: the net is a $(0,9,3)$-net over $\mathbb F_3$, and its toroidal $\delta N^{1/3}$ is $0.5291$, above the floor $0.4699$ of Theorem~\ref{thm:base23}(1) for $m\equiv0\pmod3$ ($0.3538$ for all $m$; the seed has $0.5562$).}
\newcommand{\projLinesCaption}{%
Lines in a two-dimensional projection (Remark~\ref{rem:projlines}). For odd $m=sk$, $(s,k)\neq(3,1)$, the middle digit $r^*=(m-1)/2$ of every pair of coordinates of $\mathcal P^F_{k,q}$ satisfies $y_{r^*}(x_b)\equiv\pm y_{r^*}(x_a)$ modulo more significant digits; the sign is $+1$ when one place is $\infty$, and $(-1)^{r^*}$ otherwise (checked by computation on all pairs of the net shown, and of $U^*_7$ at $q=29$, $k=1$, where $r^*=3$). This is the reversal symmetry of the two-dimensional Hammersley set (Proposition~\ref{prop:faure}), transported by M\"obius maps; $t=0$ does not see it. Shown: the frame net with the design $U^*_5$, $q=17$, $k=1$, $N=q^5$, coordinates $(x_1,x_2)$ at the places $(\infty,0)$, $r^*=2$. The relation, with balanced digits, gives the dual frequency $h=(-283,287)$, $|h|=0.338\sqrt N$, and puts $h\cdot x\bmod 1$ for all $N$ points in an interval of length at most $E=10/17$: the points lie in bands around parallel lines with normal $h$ and spacing $1/|h|=2.95\,N^{-1/2}$. (a) A window of side $12N^{-1/2}$ (light grid: $N^{-1/2}$); accent: the predicted band boundaries $h\cdot x\in\phi_0+\{0,E\}+\mathbb Z$ ($\phi_0$: start of the occupied arc); shaded: the predicted empty strips. (b) The distribution of $h\cdot x-\phi_0\bmod1$ over all $N$ points (computed): it occupies an arc of length $0.572$, and $|S(h)|=0.78$. A per-$q$ scaling of the middle generator row can break the lines without changing $t$ or the separation bounds (Remark~\ref{rem:projlines}).}
